\chapter{Criterio de traducción a claim}

\section{Traducción disciplinada a claim}
\begin{definition}
Se llama traducción a claim al paso metodológicamente controlado mediante el cual una demo seleccionada se reexpresa como afirmación verificable, con alcance declarado, dependencia visible y condiciones mínimas de validez.
\end{definition}

\begin{observation}
La traducción a claim no equivale a una proclamación fuerte por sí sola. Su función es convertir una pieza curada en una formulación más estricta, legible y contrastable.
\end{observation}

\section{Condición mínima de traducción}
\begin{definition}
Una demo seleccionada puede traducirse a claim cuando concurren simultáneamente:
\begin{enumerate}
  \item dependencia canónica visible;
  \item trazabilidad suficiente del recorrido de la demo;
  \item claridad estructural suficiente para formular una afirmación no ambigua;
  \item posibilidad de declarar alcance, condiciones de validez y condiciones de fallo;
  \item relación reconocible entre evidencia disponible y formulación propuesta.
\end{enumerate}
\end{definition}

\begin{observation}
No se exige todavía que el claim quede definitivamente validado. Se exige solo que ya pueda formularse de manera disciplinada y no meramente intuitiva.
\end{observation}

\section{Claim y demo}
\begin{definition}
Un claim derivado de una demo debe poder leerse como reformulación disciplinada de su valor estructural o metodológico, y no como expansión arbitraria de su alcance.
\end{definition}

\begin{observation}
Esta cláusula protege al Tomo III frente a sobrerreclamos que excedan la fuerza real de la pieza seleccionada.
\end{observation}

\section{Claim y evidencia}
\begin{definition}
Todo claim debe declarar, al menos de forma legible, qué parte de la evidencia disponible lo sostiene, qué parte queda abierta y qué parte solo debe leerse como hipótesis de trabajo.
\end{definition}

\begin{observation}
La traducción a claim no debe borrar la diferencia entre evidencia efectiva, interpretación disciplinada y proyección futura.
\end{observation}

\section{Condición mínima de no traducción}
\begin{definition}
Se considera que una demo no debe traducirse todavía a claim cuando su estructura siga siendo demasiado opaca, su evidencia no alcance para formular una afirmación controlada o el alcance requerido exceda la fuerza real de la pieza.
\end{definition}

\begin{observation}
Esta condición no degrada el valor de la demo. Solo evita fijar claims prematuros o inflados.
\end{observation}

\section{Cláusula de alcance}
\begin{observation}
El criterio de traducción a claim fija la frontera mínima entre demo curada y afirmación disciplinada. No constituye todavía una teoría completa de validación final, prueba fuerte o cierre exhaustivo de resultados.
\end{observation}

